\chapter{Appropriation scientifique : la dette technique et ses enjeux business}
\label{chap:veille}

% =====================================================================
%  Ce chapitre + « Travaux effectués » doivent représenter AU MOINS la
%  moitié du mémoire. C'est ici qu'on pose le cadre conceptuel qui sert
%  de grille de lecture aux 4 tâches. Mobiliser la bibliographie
%  (references.bib) : Cunningham, Fowler, Feathers, Kruchten, Humble.
% =====================================================================

\section{Définition et métaphore}
% TODO : la métaphore de Cunningham (1992) \cite{cunningham1992} ;
% « principal » vs « intérêts » ; pourquoi la dette n'est pas toujours
% une mauvaise chose (dette délibérée vs accidentelle).
\emph{[À rédiger.]}

\section{Typologie de la dette technique}
% TODO : dette de code, de test, d'architecture, de documentation,
% de dépendances/build. Quadrant de Fowler (délibérée/imprudente...).
% Relier chaque type aux tâches du mémoire.
\emph{[À rédiger.]}

\section{Enjeux business}
% TODO : le cœur du thème. Traduire la dette en termes d'entreprise :
% - vélocité de livraison et time-to-market
% - coût de maintenance et risque de régression
% - fiabilité / image (surtout pour un opérateur télécom)
% - \gls{roi} d'un investissement dans les tests / l'outillage
\emph{[À rédiger.]}

\section{Leviers de réduction}
% TODO : tests automatisés & non-régression, \gls{ci}/\gls{cd},
% refactoring \cite{fowler2018}, maîtrise du legacy \cite{feathers2004},
% documentation/outillage. Annoncer comment chaque tâche actionne un levier.
\emph{[À rédiger.]}

\section{Grille de lecture des missions}
% TODO : petit paragraphe-charnière reliant la théorie aux 4 tâches qui
% suivent (tableau optionnel : tâche -> type de dette -> levier -> bénéfice).
\emph{[À rédiger.]}
